% New analytic union review; no frozen source is changed.
\subsection{Fifteen unequal-parent eighth cells on one current tree}
\label{sec:eighth-unequal-union-review-20261009}

This review checks the analytic union, physical affine hypotheses and
accounting of the thirteen winner9 cells and the two winner50 cells.
It does not re-enumerate vertices or replay their numerical integrals.
The reviewer authored the two winner50 certificates; this note is not
described as a second-person replay of those certificates. Their full
independent numerical checks are separate bound inputs.

Write $\mu=25/128$, $\Lambda=153/256$, $\eta=10^{-6}$,
$A(s)=(1+s)/2$, $D(s)=57/64-3s/2$, $v=\min(A,D)$,
$S_j=\sum_{i\le j}t_i$, and $H_j=S_{j-1}+3t_j$.
All domains use physical ordered coordinates
$0<r=t_8<t_7<\cdots<t_1<\mu$, $t_3>1/100$ and
$r>53/13824$. The actual raw third and sixth envelopes are
\[
 L_3=\max\{v(t_1),v(t_1+t_2),v(t_1+t_3),v(t_2+t_3),
                 v(S_3),\Lambda-3t_3/8\},\qquad
 R_6=\max\{\Lambda-3t_3/8,v(s_J):\varnothing\ne J\subseteq[6]\}.
\]
The gated ordered-envelope identity and its robust finite approximation
are the already established inputs identified in the frozen
\texttt{eighth-unequal-parent-interface-20261009} interface. The new
domains and comparisons are exactly
\[
\begin{array}{c|c|c|c|c}
\text{class}&\text{number}&R=R_6&B_3&M\ \text{and selected mask}\\ \hline
9A&13&A(t_1+t_4)&D(t_2+t_3)&A(t_2+t_4+r)-\eta,\ 138\\
50A&1&A(t_2+t_5+t_6)&A(t_2+t_3)&A(t_4+t_5+t_6+r)-\eta,\ 184\\
50D&1&D(t_2+t_5+t_6)&A(t_1)&D(t_2+t_3+r)-\eta,\ 134
\end{array}
\]
Every asserted tent half has its proper affine active-side constraint.
Each domain includes $3t_1,H_3,H_5,H_7<B_3$, $M>R$,
$M-S_8<3r$, all eight initial full squares, and every remaining-selected
initial square. The independent companion checker reconstructs these
conditions from the displayed physical definitions, without importing
any author module, and checks the complete affine-row multisets and
strict rational centers of all fifteen domains.

\paragraph{The baseline minorant in 50D: an explicit extension.}
The frozen interface's displayed (U1) literally specifies
$B_3=v(t_2+t_3)$. It must not be cited as a literal statement about
$A(t_1)$. Here is the required specialization of its proof.
On $t_1<\mu$ we have $A(t_1)=v(t_1)$, one of the actual $L_3$
entries; hence $B_3=A(t_1)\le L_3$. On a compact subset of the
union of the fifteen strict domains let
\[
 \delta=\min\{B_3(t)-H_j(t):j=3,5,7\}>0,
\]
where each component uses its own displayed minorant. On the finite
union of first-three-coordinate projections, approximate the same
true $L_3$ from below with a common causal coarse rule
$L_3-\delta/3<\lambda_3\le L_3$. This is permitted by the existing
uniform compact approximation. Then all three histories are below
$\lambda_3$ in every component. The choice depends only on the
first-three labels, never separately on an eighth child. This argument
uses only $B_3\le L_3$ and the strict histories, so it applies also to
the two pair minorants already stated in the frozen interface.

Thus no new third distribution level, pair license at $A(t_1)$, or
assertion $\lambda_3=B_3$ is needed. The other third branches may be
larger than $B_3$. This new paragraph proves the baseline instance
explicitly while leaving the older, more narrowly worded interface
unchanged. It is a common sufficiently fine finite realization, not
a claim that every earlier coarse realization already has these paths.

\paragraph{All selected gates and the actual parent.}
For a selected mask $I$, let $b=s_I$ and enumerate its complement
in increasing slot order as $q_1,\ldots,q_m$. Besides
$S_{j-1}+2t_j<M$ for every $1\le j\le8$, the necessary initial tests are
\[
 b+\sum_{a<k}t_{q_a}+2t_{q_k}<M\quad(1\le k\le m).
\]
The complement lists checked by the companion are respectively
$(1,3,5,6,7)$ for138, $(1,2,3,7)$ for184, and $(1,4,5,6,7)$
for134. In particular a stale ``selected184'' name in the descending
source labels does not alter its coefficients: they are exactly the
five mask134 tests. The active cap identity and the other tent-half
inequality imply $M=v(b)-\eta$ throughout each component.
Once the eighth slot has appeared, every later selected remaining
test is the full square test. Ordinary upper residual odd cubes
give the subsequent full square support. These are precisely the
established selected-product factorization hypotheses, including
the capacity-below-one repair and the signed fixed-label compilation;
no atom is replaced by an unjustified positive arithmetic weight.

Keep every already adopted fourth-, fifth- and sixth-level action on
the common refinement. Their actual cap satisfies
\[
 \lambda_3\le\lambda_6^{\rm current}\le R_6=R.
\]
The upper inequality follows because every sixth increase is below
a first-six selected tent, and inherited caps are below the same
all-subset envelope. A comparison-envelope or moment improvement
does not change a cap. In particular the four adopted nonsingleton
sixth cells must not be discarded merely because they can overlap
the present cells. Their labels and gains remain in the current tree.

The seventh history is retained since $H_7<\lambda_3$. More strongly,
\[
 \lambda_6^{\rm current}-S_8>H_7-S_8=2t_7-r>r,
\]
so
\[
 1<\frac{\lambda_6^{\rm current}-S_8}{r}
 \le\frac{R-S_8}{r}<\frac{M-S_8}{r}<3.
\]
This proves the ordinary reciprocal range for the actual old residual,
not only for its upper comparison. At an even eighth node the upper
residual's odd first exits have odd full rank. Prefix closure and upper
majorization survive. Its exact raw density difference is the old
ordinary residual density minus the new one. Because
$\bar F(s)=1/s$ in this range, it is bounded below at the endpoint by
\[
 \frac{4(M-R)}{(R-S_8)(M-S_8)\prod_{i=1}^8t_i}.
 \tag{J}
\]
Increasing the old comparison cap to $R$ decreases this difference.
There is no assumption $\lambda_6^{\rm current}=R$, and no monotonicity
claim about individual finite signed coefficients. Nor is a Chen
quarter or an ordering factorial applied to this physical upper-$K$
density. On an eight-simplex, the affine gap mean is the sum of its
nine vertex gaps divided by9, and each gap-weighted affine factor
mean has denominator10; this is the same bound interface as before.

\paragraph{The complete disjointness argument.}
All 63 old chosen planes are at most $R$, and the asserted winning
tent actually equals $R$. Every first-three candidate's chosen plane
and $W=\Lambda-3t_3/8$ is a proper, nonconstant constrained affine
function strictly below $R$ in the interior. Therefore the true
envelope, not merely its minorant, obeys
$L_3<R_6$ on all fifteen interiors. All thirty earlier adopted eighth
domains obey $R_6=L_3$, so none acts on these new interiors.

For the thirteen9A components put $e=t_1+t_4<\mu$.
For any first-six subset sum the two allowed old-plane conditions are
$s\le e$ and $s\ge T=25/96-e/3$, with
$T-e=4(\mu-e)/3>0$. Any two different saved complete patterns differ
on a subset, so their interiors cannot meet. All 78 such pairs are
checked. Family cell11 has exactly the seed's affine-row multiset;
it is excluded from the twelve-new-cell sum and the seed is counted
once. The two rejected floating navigation nodes are not used as
infeasibility proofs: no exhaustive unequal-parent classification
is asserted here.

In either50 component the mask9 chosen-plane bound is proper and
strict, so $v(t_1+t_4)<R_6$. In any9A component that same tent equals
$R_6$. This separates all26 cross pairs without relying on the names
of the winners. Finally the50A half has $t_2+t_5+t_6<\mu$ while50D
has that sum $>\mu$. Together these give all105 interior pair
separations. Proper equality walls are removed in a compact exhaustion
or assigned deterministically; none receives two positive-volume
integrals. The winning and active-cap identities are retained, not
trimmed as if they were proper walls.

Possible overlap with sixth-level domains is instead handled by
\[
 V_{\rm base}-V_{\rm next}
 =(V_{\rm base}-V_{\rm current})
  +(V_{\rm current}-V_{\rm next}).
\]
Formula (J) bounds only the second, eighth-level difference. It does
not add another comparison from the obsolete fifth cap or re-add a
sixth gain. The old seventh decision and all earlier histories remain
attached to the original parent labels.

\paragraph{Finite construction, accounting and audit scope.}
All denominators are positive on the closed cells: order and the
uniform positive $r$ cutoff bound the coordinate factors, while
the history gives $R-S_8\ge r$ and $M-S_8\ge r$. Thus proper boundary
strips have vanishing integral. Trim a finite union of such strips,
freeze the common actual third cap first, refine and retain current
ancestor labels, then freeze new eighth caps below their licenses
and above all tests and inherited caps. Revert to the inherited rule
on unused labels; do not introduce even first exits. All buffers and
partitions are fixed before taking $N\to\infty$. Constants may depend
on this fixed construction. Existing approximation/compiler budgets
must be retained and new losses paid separately.

The author ledger to be used once, subject to separate numerical
verification, is
\[
 \Delta_{15}=\Delta_{9,\rm seed}+\Delta_{9,\rm new12}
                 +\Delta_{50A}+\Delta_{50D}.
\]
No narrow seed model or repeated family cell11 is included. The
companion exact checker reconstructs the physical rows, strict centers,
seed identity and separation ledger and checks the rational addition;
it does not enumerate any bases, triangulations or Jensen floors.
The frozen independent50 numerical reports are recorded separately.
This review supplies the local analytic union and the baseline-minorant
extension only. It finds no missing ancestry or disjointness condition
within the stated fifteen-cell class. It does not reprove external
distribution theorems, claim all eighth strategies are covered, or
select a new global $K,c,b$.
